\documentclass[reqno, 12pt]{amsart}

\usepackage{amssymb}
\usepackage{amsthm}
\usepackage{amsmath}
\usepackage{amsxtra}
\usepackage{mathrsfs}
\usepackage{comment}
\usepackage{graphicx}
\usepackage{placeins}
\usepackage{multicol}
\usepackage{bbm}
\usepackage{stmaryrd}
%\usepackage{enumerate}
\usepackage[inline, shortlabels]{enumitem}
\usepackage{mathdots}
\usepackage{colonequals}
\usepackage{hyperref}


\usepackage{calc}
\usepackage{tikz}
\usetikzlibrary{arrows,calc,automata,shadows,backgrounds,positioning,intersections,fadings,decorations.pathreplacing,shapes,matrix}
\usepackage{tikz-cd}

\usepackage{fullpage}

\newtheorem{thm}{Theorem}[section]
\newtheorem{lem}[thm]{Lemma}
\newtheorem{cor}[thm]{Corollary}
\newtheorem{conj}[thm]{Conjecture}
\newtheorem{prop}[thm]{Proposition}
\theoremstyle{definition}  
\newtheorem{defn} [thm] {Definition} 
\newtheorem{claim}[thm]{Claim}
\newtheorem{example} [thm] {Example}
\newtheorem{rem} [thm] {Remark}

\newenvironment{problab}[1]
{\noindent\textbf{Problem #1}.}
{\vskip 6pt}

\newenvironment{exolab}[1]
{\noindent\textbf{Exercise #1}.}
{\vskip 6pt}

\theoremstyle{remark}
\newtheorem*{soln}{Solution}

\newcommand{\C}{\mathbb C}
\renewcommand{\H}{\mathbb H}
\newcommand{\D}{\mathbb D}
\newcommand{\F}{\mathbb F}
\newcommand{\Q}{\mathbb Q}
\newcommand{\R}{\mathbb R}
\newcommand{\Z}{\mathbb Z}
\newcommand{\T}{\mathbb T}
\renewcommand{\P}{\mathcal P}
\renewcommand{\O}{\mathcal O}
\newcommand{\m}{\mathfrak m}
\newcommand{\A}{\mathbb A}
\renewcommand{\SS}{\mathbb S}
\renewcommand{\L}{\mathcal L}

\newcommand{\B}{\mathcal B}
\newcommand{\E}{\mathcal E}

\newcommand{\SSpan}{\text{span}}

\DeclareMathOperator{\lcm}{lcm}
\DeclareMathOperator{\img}{img}
\DeclareMathOperator{\Ann}{Ann}
\DeclareMathOperator{\Tor}{Tor}
\DeclareMathOperator{\tr}{tr}
\DeclareMathOperator{\Hom}{Hom}
\DeclareMathOperator{\End}{End}
\DeclareMathOperator{\Aut}{Aut}
\DeclareMathOperator{\Gal}{Gal}
\DeclareMathOperator{\sgn}{sgn}
\DeclareMathOperator{\ord}{ord}
\DeclareMathOperator{\ad}{ad}
\DeclareMathOperator{\coker}{coker}
\DeclareMathOperator{\rank}{rank}

\DeclareMathOperator{\id}{id}

\usepackage{mathpazo}

\everymath{\displaystyle}

\tikzset{commutative diagrams/.cd, arrow style = tikz, diagrams = {>=latex}}

\newenvironment{amatrix}[1]{%
  \left(\begin{array}{@{}*{#1}{c}|c@{}}
}{%
  \end{array}\right)
}

\usepackage{abstract}
\renewcommand{\abstractname}{}    % clear the title
\renewcommand{\absnamepos}{empty} % originally center
%\setlength{\abstitleskip}

\begin{document}
%\renewcommand{\abstractname}{\vspace{-\baselineskip}}

\title{18.700 Problem Set 4}
%\author{Évariste Galois}
\dedicatory{Due Wednesday, October 2 at 11:59 pm on \href{https://canvas.mit.edu/courses/27315}{Canvas}}
%\thanks{}
%\contrib{}
%\address{}

\maketitle

\begin{abstract}
  \noindent Collaborated with:\\
  Sources used: 
\end{abstract}
\vskip 0.5cm

\begin{problab}{1} (3A \#7) (3 points)
  Prove that if $V$ is a $1$-dimensional vector space and $T \in \L(V)$, then there exists a scalar $\lambda \in \F$ such that $T(v) = \lambda v$ for all $v \in V$.
\end{problab}

%\begin{soln}
%\end{soln}

\begin{problab}{2} (8 points)
  \begin{enumerate}[(a)]
    \item 
      Prove that there exists a linear map $T: \F^2 \to \F^3$ such that
      \begin{equation*}
        T \begin{pmatrix}1\\ 1\end{pmatrix} = \begin{pmatrix}1\\ 0\\ 2\end{pmatrix}
        \qquad
        \text{and}
        \qquad
        T \begin{pmatrix}2\\ 3\end{pmatrix} = \begin{pmatrix}1\\ -1\\ 4\end{pmatrix} \, .
      \end{equation*}
    \item
      What is $T\left(\begin{smallmatrix}x\\ y\end{smallmatrix}\right)$ for any $x,y \in \F$?
    \item
      Is $T$ one-to-one?
  \end{enumerate}
\end{problab}

\begin{problab}{3} (2C \#4) (8 points)
  \begin{enumerate}[(a)]
    \item
      Let $U \colonequals \{p \in \P_4(\R) : p''(6) = 0\}$. Find a basis of $U$.
    \item
      Extend the basis in (a) to a basis of $\P_4(\R)$.
    \item
      Find a subspace $W$ of $\P_4(\R)$ such that $\P_4(\R) = U \oplus W$.
  \end{enumerate}
\end{problab}

\begin{problab}{4} (3B \#3) (8 points)
  Suppose $v_1, \ldots, v_m$ is a list of vectors in $V$. Define $T \in \L(\F^m, V)$ by
  \begin{equation*}
    T(z_1, \ldots, z_m) \colonequals z_1 v_1 + \cdots + z_m v_m \, .
  \end{equation*}
  \begin{enumerate}[(a)]
    \item
      What property of $T$ corresponds to $v_1, \ldots, v_m$ spanning $V$?
    \item
      What property of $T$ corresponds to the list $v_1, \ldots, v_m$ being linearly independent?
  \end{enumerate}
  Give proofs of your claims in both cases. (Be sure to prove both directions of any ``iff'' statements.)
\end{problab}

\begin{problab}{5} (2A \#3, \#14, 2B \#9) (10 points)
  Suppose $v_1, \ldots, v_m$ is a list of vectors in $V$. For $k \in \{1, \ldots, m\}$, let
  \begin{equation*}
    w_k \colonequals v_1 + \cdots + v_k \, .
  \end{equation*}
    \begin{enumerate}[(a)]
      \item 
        Show that $\SSpan(v_1, \ldots, v_m) = \SSpan(w_1, \ldots, w_m)$.
      \item
        Show that the list $v_1, \ldots, v_m$ is linearly independent iff the list $w_1, \ldots, w_m$ is linearly independent.
      \item
        Show that $v_1, \ldots, v_m$ is a basis of $V$ iff $w_1, \ldots, w_m$ is a basis of $V$.
    \end{enumerate}
    (\textit{Hint:} It may be useful to find a formula for $v_k$ in terms of $w_1, \ldots, w_m$.)
\end{problab}

\end{document} 	
